% Partie 4 du notebook : repérer les incidents.

\bandeaupartie{I. Partie 4 : repérer les incidents}{Plus de plaintes qu'attendu, sur une seule marque : c'est la signature d'un incident.}

\section{Partie 4 : repérer les incidents}

\subsection{Observé contre attendu}

Prenons un problème $g$ (par exemple \og goût de savon \fg) et une case : une marque $b$ et une période $t$.
Si ce problème se répartissait au hasard entre les cases, on en attendrait
\begin{formulabox}
\[
\lambda_{g,b,t} = \frac{n_g \times n_{b,t}}{N},
\]
où $n_g$ est le nombre de plaintes du problème $g$, $n_{b,t}$ le nombre de plaintes (tous problèmes
confondus) de la marque $b$ à la période $t$, et $N$ le nombre total de plaintes. Le notebook note cette
valeur \texttt{attendu} ; on la compare au nombre observé $o_{g,b,t}$.
\end{formulabox}

\begin{center}
\begin{tikzpicture}[x=0.95cm, y=0.62cm, font=\footnotesize]
\node[font=\bfseries\small, color=darkblue] at (5,6.3) {problème \og goût de savon \fg{} : plaintes observées (attendues)};
\foreach \p [count=\j] in {26 P03,26 P04,26 P05,26 P06,26 P07} \node[color=slate] at (\j*1.9-0.3,5.4) {\p};
\foreach \m [count=\i] in {Brand 001,Brand 002,Brand 003,Brand 004} \node[color=slate, anchor=east] at (0.4,5.4-\i) {\m};
\foreach \i/\j/\o/\e in {1/1/1/{0,9},1/2/0/{0,8},1/3/1/{1,0},1/4/0/{0,9},1/5/1/{0,8},
  2/1/0/{0,6},2/2/1/{0,7},2/3/0/{0,7},2/4/1/{0,6},2/5/0/{0,6},
  3/1/1/{0,8},3/2/0/{0,9},3/4/1/{0,9},3/5/0/{0,8},
  4/1/1/{1,0},4/2/1/{0,9},4/3/1/{1,1},4/4/0/{0,9},4/5/1/{1,0}} {
  \fill[notebg] (\j*1.9-1.15,5.4-\i-0.42) rectangle (\j*1.9+0.55,5.4-\i+0.42);
  \node at (\j*1.9-0.3,5.4-\i) {\o\ (\e)};
}
\fill[accentblue] (3*1.9-1.15,5.4-3-0.42) rectangle (3*1.9+0.55,5.4-3+0.42);
\node[white, font=\footnotesize\bfseries] at (3*1.9-0.3,5.4-3) {9 (2,0)};
\draw[-{Stealth}, accentblue, thick] (11,1.9) node[right, text width=3.3cm, font=\footnotesize] {9 plaintes là où le hasard en donnerait environ 2} -- (6.4,2.3);
\end{tikzpicture}
\end{center}
{\small\color{slate}\textit{Tableau inventé, pour illustrer le calcul.}}

\subsection{La loi de Poisson : est-ce le hasard ?}

\begin{formulabox}
Si le hasard seul jouait, le nombre de plaintes $X$ de la case suivrait une loi de Poisson de paramètre
$\lambda = \lambda_{g,b,t}$. La \textbf{p-valeur} est la probabilité d'observer au moins autant :
\[
p = P(X \geq o) = 1 - \sum_{k=0}^{o-1} \frac{\lambda^{k}\, \mathrm{e}^{-\lambda}}{k!},
\]
où $o = o_{g,b,t}$ est le nombre observé, $k$ parcourt les valeurs entières de $0$ à $o-1$, et
$\mathrm{e} \approx 2{,}718$ est la base du logarithme népérien.
\end{formulabox}
\illus[0.82\linewidth]{ill_poisson.png}{Exemple inventé : avec 2 plaintes attendues, en observer 9 arrive environ 2 fois sur 10\,000.}

\subsection{Tester des centaines de cases : la correction de Benjamini-Hochberg}

\begin{tcbraster}[raster columns=2, raster equal height, raster column skip=8pt, raster every box/.style={unbreakable}]
\begin{enfant}[unbreakable]
\textbf{Lancer cent pièces.} Si on lance cent pièces dix fois, l'une d'elles finira bien par faire
\og pile \fg{} huit fois : ce n'est pas une pièce truquée, c'est qu'on a regardé beaucoup de pièces.
Tester des centaines de cases, c'est pareil.
\end{enfant}
\begin{formulabox}
On range les $m$ p-valeurs par ordre croissant, $p_{(1)} \leq \ldots \leq p_{(m)}$, puis
\[
q_{(i)} = \min_{j \geq i}\, \min\!\left(1,\ \frac{m\, p_{(j)}}{j}\right).
\]
Garder les cases avec $q < 0{,}01$ limite à environ 1 \% la part de fausses découvertes parmi elles.
\end{formulabox}
\end{tcbraster}

\subsection{Une marque, ou tout le monde ? La part dans le pic}

\begin{formulabox}
\[
\pi_{g,b,t} = \frac{o_{g,b,t}}{o_{g,t}},
\]
où $o_{g,t}$ est le nombre de plaintes du problème $g$ à la période $t$, toutes marques confondues. Un
incident de production touche une marque : $\pi$ est grand. Une saison touche tout le monde : $\pi$ est petit.
\end{formulabox}

\begin{stepbox}[title={4.1\quad Les cases anormales}]
\texttt{cases\_anormales(cles, seuil=0.01, o\_min=5)} calcule, pour chaque problème et chaque case :
observé, attendu, p-valeur, $q$. Elle garde les cases avec $q < 0{,}01$ et au moins 5 plaintes.
On l'applique à (marque, période) en ajoutant \texttt{part\_pic}, puis à la période seule.
\doisvoir{des cases anormales sur marque et période ; toutes les valeurs $q$ et toutes les parts entre 0 et 1.}
\end{stepbox}

\begin{enfant}[unbreakable]
\textbf{Le sac de bonbons.} Un sac où 10 \% des bonbons sont rouges : en en tirant 20, on attend environ
\textbf{2 rouges}, c'est $\lambda$. On en tire \textbf{9} : la \textbf{p-valeur} est la chance d'en avoir
au moins autant par pure malchance ; si elle est minuscule, le sac a un problème. Mais on regarde des
centaines de sacs, et l'un d'eux finira bien par sortir beaucoup de rouges par hasard : \textbf{$q$} est la
p-valeur rendue plus sévère pour en tenir compte. On ne lève une alerte que si $q < 0{,}01$.
\end{enfant}

\begin{stepbox}[title={4.2\quad Incident, pic partagé ou campagne ? Dans cet ordre}]
\begin{center}
\begin{tikzpicture}[
  q/.style={draw=darkblue, fill=notebg, rounded corners=4pt, text width=4.6cm, align=center, font=\footnotesize, minimum height=1.05cm},
  r/.style={draw=#1, fill=#1, text=white, rounded corners=4pt, text width=2.6cm, align=center, font=\footnotesize\bfseries, minimum height=0.9cm},
  fl/.style={-{Stealth[length=2.5mm]}, thick, color=slate}]
\node[q] (q1) {plus de la moitié de ses plaintes ont une voisine au-delà de 0,9 ?};
\node[q, below=0.75cm of q1] (q2) {une case anormale où la marque porte au moins la moitié du pic ($\pi \geq 0{,}5$) ?};
\node[q, below=0.75cm of q2] (q3) {une case ou une période anormale, sans marque qui porte le pic ?};
\node[r=cacao, right=1.4cm of q1] (r1) {campagne};
\node[r=accentblue, right=1.4cm of q2] (r2) {incident candidat};
\node[r=ambre, right=1.4cm of q3] (r3) {pic partagé};
\node[r=slate, below=0.75cm of q3] (r4) {fond};
\draw[fl] (q1) -- node[above, font=\footnotesize] {oui} (r1);
\draw[fl] (q2) -- node[above, font=\footnotesize] {oui} (r2);
\draw[fl] (q3) -- node[above, font=\footnotesize] {oui} (r3);
\draw[fl] (q1) -- node[right, font=\footnotesize] {non} (q2);
\draw[fl] (q2) -- node[right, font=\footnotesize] {non} (q3);
\draw[fl] (q3) -- node[right, font=\footnotesize] {non} (r4);
\end{tikzpicture}
\end{center}
Les cases des incidents candidats forment les \texttt{alertes}.
\doisvoir{chaque problème a un statut : incident candidat, campagne, pic partagé ou fond. Une base peut très bien n'avoir aucun incident.}
\end{stepbox}

\begin{enfant}
\textbf{À la cantine de l'école.} Les classes sont les marques, les semaines sont les périodes.
\end{enfant}
\begin{tcbraster}[raster columns=2, raster equal height, raster column skip=6pt, raster row skip=6pt, raster every box/.style={unbreakable}]
\begin{kpibox}[title={\faBullhorn\ Campagne}, unbreakable]
\small 30 élèves recopient la même lettre : \og on veut des frites \fg. C'est une pétition, pas un problème
de cuisine.
\end{kpibox}
\begin{kpibox}[title={\faExclamationTriangle\ Incident candidat}, unbreakable]
\small Mardi, 8 élèves du CM1, et presque personne d'autre, ont mal au ventre après le yaourt. Un mauvais
lot ?
\end{kpibox}
\begin{kpibox}[title={\faSun\ Pic partagé}, unbreakable]
\small En décembre, toutes les classes toussent en même temps. Ce n'est pas la cantine : c'est l'hiver.
\end{kpibox}
\begin{kpibox}[title={\faMugHot\ Fond}, unbreakable]
\small Chaque semaine, un ou deux élèves trouvent la soupe trop chaude. Rien d'anormal : c'est la vie
normale.
\end{kpibox}
\end{tcbraster}

\illus[0.98\linewidth]{4_carte_chaleur_problemes_periodes.png}{Figure 4 du notebook (base fictive) : chaque problème, période par période.}

\begin{lire}
\textbf{Comment lire la carte de chaleur.} Une ligne par problème, une colonne par période. La couleur
compare les plaintes observées à celles attendues si le problème était réparti au hasard dans le temps :
crème quand c'est comme attendu, rouge de plus en plus foncé pour deux, quatre, huit fois plus, taupe pâle
pour moins ou aucune. Une ligne rouge sur une ou deux cases seulement : un problème qui surgit puis
disparaît. Les cases encadrées sont les \textbf{alertes} ; la couleur du nom de la ligne donne le statut
(rouge : incident, ambre : pic partagé, cacao : campagne).
\end{lire}
